\section{Proposta: o padrão da indistinguibilidade sob a LGPD}

% Objetivo da seção: formular o padrão técnico-jurídico em linguagem da LGPD
% e mostrar que ele decorre dos arts. 12, 16 e 18, IV; o RGPD entra só para
% demonstrar compatibilidade, não como fundamento.

Esta seção deve enunciar o padrão: a eliminação de dados pessoais de um
modelo está satisfeita quando a contribuição do titular se torna
indistinguível da de um modelo que nunca viu esses dados, avaliada contra o
adversário que o art.~12 da LGPD pressupõe (meios técnicos razoáveis e
disponíveis na ocasião). Deve mostrar que esse resultado é, juridicamente,
uma anonimização, e portanto um dos remédios que o art.~18, IV já coloca ao
lado da eliminação; e que o art.~16 fornece a válvula de escape quando nem
a indistinguibilidade for tecnicamente viável, mediante conservação
documentada e com finalidade limitada. Deve especificar os componentes
operacionais: modelo de adversário, métrica de verificação, ônus da prova
do controlador e documentação exigível pela ANPD. Ao final, deve indicar
brevemente que o mesmo padrão é compatível com o Considerando~26 do RGPD e
com EDPS~v~SRB, embora o RGPD careça de cláusula equivalente ao art.~16.

\todo{Formulação em linguagem normativa (proposta de enunciado interpretativo do art.~18, IV).}
\todo{Como escolher o adversário à luz do art.~12: capacidades, acesso ao modelo, dados auxiliares.}
\todo{Relação com os três testes do Parecer 28/2024 como parâmetro de adversário.}
\todo{Ônus da prova (art.~6, X, responsabilização e prestação de contas) e papel de auditoria independente.}
